\documentclass[10pt,a4paper]{article}

\usepackage[margin=0.58in]{geometry}
\usepackage[T1]{fontenc}
\usepackage[utf8]{inputenc}
\usepackage{lmodern}
\usepackage{tabularx}
\usepackage{enumitem}
\usepackage{xcolor}
\usepackage{fontawesome5}
\usepackage{hyperref}

\definecolor{accent}{HTML}{1F4E79}

\hypersetup{
  colorlinks=true,
  linkcolor=accent,
  urlcolor=accent,
  citecolor=accent,
  pdfborder={0 0 0}
}

\pagestyle{empty}
\setlength{\parindent}{0pt}
\setlength{\parskip}{0pt}
\setlength{\emergencystretch}{2em}
\renewcommand{\baselinestretch}{1.02}

\setlist[itemize]{
  leftmargin=1.15em,
  itemsep=0.18em,
  topsep=0.2em,
  parsep=0pt,
  partopsep=0pt
}

% Section title with a light rule below it.
\newcommand{\sectionTitle}[1]{%
  \vspace{8pt}
  \par\noindent{\normalsize\bfseries\textcolor{accent}{\MakeUppercase{#1}}}\par
  \vspace{3pt}
  \hrule height 0.45pt
  \vspace{4.5pt}
}

% Project title on the left, repository label/link on the right, bullets below.
\newcommand{\resumeProject}[3]{%
  \vspace{3.5pt}
  \begin{tabularx}{\textwidth}{@{}>{\raggedright\arraybackslash}X@{\hspace{1.2em}}r@{}}
    \textbf{#1} & #2 \\
  \end{tabularx}
  \vspace{-0.55em}
  #3
}

\begin{document}
\raggedright

\begin{center}
  {\LARGE \textbf{\textcolor{accent}{Bui Ngoc Thai}}}\\
  \vspace{0.18em}
  DevOps Intern \textbar{} Automation, CI/CD \& Cloud Infrastructure\\
  \vspace{0.25em}
  {\footnotesize
  \faMapMarker*~Ho Chi Minh City, Vietnam \textbar{}
  \faEnvelope~\href{mailto:23521412@gm.uit.edu.vn}{23521412@gm.uit.edu.vn} \textbar{}
  \faPhone*~(+84) 346 072 516 \textbar{}
  \faGithub~\href{https://github.com/Nashitaz13}{github.com/Nashitaz13} \textbar{}
  \faLinkedin~\href{https://www.linkedin.com/in/nashitaz13/}{linkedin.com/in/nashitaz13}}
\end{center}

\vspace{-0.55em}

\sectionTitle{Summary}
Final-year Computer Networks and Data Communications student with an automation-first mindset and interest in DevOps, CI/CD, and cloud infrastructure. I enjoy turning manual setup and deployment steps into repeatable scripts, pipelines, and documented workflows while troubleshooting issues across Linux, networking, and Kubernetes-based environments.

\sectionTitle{Education}
\begin{tabularx}{\textwidth}{@{}>{\raggedright\arraybackslash}Xr@{}}
  \textbf{University of Information Technology, VNU-HCM} & Expected Graduation: 2027 \\
\end{tabularx}
Major in Computer Networks and Data Communications \textbar{} Final-year student \textbar{} GPA: 8.33/10

\sectionTitle{Projects}
\resumeProject{Drift-Aware MLOps Framework for Multi-Model Serving --- Research \& Prototype}{%
  \footnotesize Under Review \textbar{} LNCS/Springer}{%
\begin{itemize}
  \item Co-authored a research paper formalising a drift-aware MLOps lifecycle with a tenant-aware
        control-plane/data-plane split, a Native Model Registry backed by PostgreSQL, and an adaptation
        policy $\pi$ that maps feature-level drift scores to review-driven retraining triggers.
  \item Implemented the prototype on K3s/AWS: Django/DRF control plane, FastAPI dynamic prediction
        endpoints, asynchronous Redpanda log streaming, Evidently AI drift reports, Argo Workflows
        + Kubeflow \texttt{PyTorchJob} training orchestration, and Karpenter node autoscaling.
  \item Validated via offline multi-level replay on CIC-IDS2017 across 9 scenarios (3 attack
        families $\times$ 3 severity levels), repeated over 5 random seeds (std $\leq$ 0.004).
        Post-hoc Bonferroni and BH corrections verified no drift decision changed.
        Retrained challengers recovered $\Delta$Macro-F1 by +0.54 to +0.87 depending on severity.
\end{itemize}
}

\resumeProject{NT531 EKS Cilium vs kube-proxy Benchmark}{\footnotesize\href{https://github.com/daithang59/NT531_EKS-Cilium-Kubeproxy-Benchmark}{GitHub: daithang59/NT531\_EKS-Cilium-Kubeproxy-Benchmark}}{
\begin{itemize}
  \item Evaluated Kubernetes Service datapath performance across 2 modes, 3 scenarios, and 3 calibrated load levels by running Fortio benchmarks on AWS EKS.
  \item Analyzed p99 latency, RPS, error rate, 95\% confidence intervals, and repeatability CV\%, supported by Prometheus/Grafana metrics and Hubble flow evidence.
\end{itemize}
}

\resumeProject{SageLMS -- DevSecOps CI/CD \& Microservices Learning Platform}{\footnotesize\href{https://github.com/daithang59/sagelms}{GitHub: daithang59/sagelms}}{
\begin{itemize}
  \item Supported DevSecOps CI/CD workflows for a microservices-based SageLMS platform, including PR validation, Docker image builds, Trivy scans, SBOM generation, Harbor push, and GitOps deploy PRs.
  \item Documented and validated container supply-chain evidence, including image tags/digests, Trivy reports, CycloneDX SBOMs, Harbor pull secret flow, and Cosign verification logs.
\end{itemize}
}

\resumeProject{MentorMe Mobile App -- Mobile \& Backend Integration Project}{\footnotesize\href{https://github.com/kaing615/MentorMe-Mobile-App}{GitHub: kaing615/MentorMe-Mobile-App}}{
\begin{itemize}
  \item Contributed to a full-stack mentoring platform across 3 components: Android Jetpack Compose app, Express/TypeScript backend, and React Admin dashboard.
  \item Supported team development across Android, backend, and admin modules using Gradle/npm build scripts, TypeScript checks, Swagger API docs, Android test structure, and PR-based Git workflow.
\end{itemize}
}

\sectionTitle{Technical Skills}
\begin{itemize}
  \item \textbf{Programming/Scripting:} Python, Bash, PowerShell, Java, TypeScript, Kotlin, C/C++
  \item \textbf{DevOps \& CI/CD:} Git, GitHub Actions, Jenkins, Docker, Docker Compose, Kubernetes/K3s, Argo CD
  \item \textbf{Cloud \& Infrastructure:} AWS EC2/S3/IAM/OIDC, Terraform, Helm, Linux, Nginx, Cloudflare
  \item \textbf{Monitoring \& Testing:} Prometheus, Grafana, API smoke testing, Locust, integration testing, benchmark automation
  \item \textbf{Networking \& Databases:} OSI model, DNS, DHCP, firewall concepts, PostgreSQL, MongoDB, Redis
\end{itemize}

\end{document}
